% google-fonts.tex — Compatible pdfLaTeX / XeLaTeX / LuaLaTeX

% 1) Détecter le moteur
\usepackage{iftex}

% 2) Couleurs et titres (utilisés plus bas)
\usepackage{xcolor}
\usepackage{titlesec}

% 3) Fonts selon le moteur
\ifPDFTeX
  % ----- pdfLaTeX : pas de Google Fonts -----
  \usepackage[T1]{fontenc}
  \usepackage[utf8]{inputenc}
  \usepackage{lmodern} % ou newtxtext,newtxmath si tu préfères
\else
  % ----- XeLaTeX / LuaLaTeX : Google Fonts via fontspec -----
  \usepackage{fontspec}
  \defaultfontfeatures{Ligatures=TeX, Scale=MatchLowercase}

  % --- Roboto en texte et sans-serif ---
  % Place les fichiers dans: fonts/Roboto/Roboto-*.ttf
  \setmainfont{Roboto}[
    Path=fonts/Roboto/,
    Extension=.ttf,
    UprightFont   = *-Regular,
    BoldFont      = *-Bold,
    ItalicFont    = *-Italic,
    BoldItalicFont= *-BoldItalic
  ]
  \setsansfont{Roboto}[
    Path=fonts/Roboto/,
    Extension=.ttf,
    UprightFont   = *-Regular,
    BoldFont      = *-Bold,
    ItalicFont    = *-Italic,
    BoldItalicFont= *-BoldItalic
  ]

  % --- (Optionnel) Roboto Mono pour le code ---
  % Fichiers dans: fonts/RobotoMono/RobotoMono-*.ttf
  % \setmonofont{RobotoMono}[
  %   Path=fonts/RobotoMono/,
  %   Extension=.ttf,
  %   UprightFont   = *-Regular,
  %   BoldFont      = *-Bold,
  %   ItalicFont    = *-Italic,
  %   BoldItalicFont= *-BoldItalic
  % ]
\fi

% Amélioration des sections avec Google Fonts (sans toucher aux chapitres)
\titleformat{\section}
  {\fontsize{16pt}{19pt}\selectfont\bfseries\color{blue!80!black}\sffamily}
  {\thesection}{1em}{}

\titleformat{\subsection}
  {\fontsize{14pt}{17pt}\selectfont\bfseries\color{blue!70!black}\sffamily}
  {\thesubsection}{1em}{}

\titleformat{\subsubsection}
  {\fontsize{12pt}{15pt}\selectfont\bfseries\color{blue!60!black}\sffamily}
  {\thesubsubsection}{1em}{}

% Amélioration des paragraphes et du texte général
\ifPDFTeX
  % Pour pdfLaTeX, améliorer l'espacement et la lisibilité
  \renewcommand{\baselinestretch}{1.15}
  \setlength{\parskip}{6pt plus 2pt minus 1pt}
  \setlength{\parindent}{1.5em}
  
  % Améliorer la typographie avec microtype
  \usepackage[activate={true,nocompatibility},final,tracking=true,kerning=true,spacing=true,factor=1100,stretch=10,shrink=10]{microtype}
  
  % Configuration pour de meilleures polices en pdfLaTeX
  \usepackage[scaled=0.92]{helvet} % Pour une meilleure police sans-serif
  % Garder la police par défaut (serif) pour le contenu
\else
  % Pour XeLaTeX/LuaLaTeX avec Google Fonts
  \renewcommand{\baselinestretch}{1.15}
  \setlength{\parskip}{6pt plus 2pt minus 1pt}
  \setlength{\parindent}{1.5em}
  
  % Améliorer l'affichage des citations et éléments spéciaux
  \newfontfamily\quotefont{Roboto}[
    Path=fonts/,
    Extension=.ttf,
    UprightFont=*-Italic
  ]
\fi

% Style pour les listes avec de meilleurs espacements
\usepackage{enumitem}
\setlist[itemize]{
  leftmargin=1.5em,
  itemsep=4pt plus 2pt minus 1pt,
  parsep=2pt plus 1pt minus 1pt,
  topsep=6pt plus 2pt minus 2pt
}

\setlist[enumerate]{
  leftmargin=1.8em,
  itemsep=4pt plus 2pt minus 1pt,
  parsep=2pt plus 1pt minus 1pt,
  topsep=6pt plus 2pt minus 2pt
}

% Amélioration des légendes de figures et tableaux
\usepackage{caption}
\captionsetup{
  font={small,sf},
  labelfont={bf,color=blue!80!black},
  textfont={color=black},
  margin=10pt,
  skip=8pt
}

% Style pour les références bibliographiques
\usepackage{url}
\renewcommand{\UrlFont}{\small\ttfamily}

% Amélioration de l'affichage des footnotes
\renewcommand{\footnotesize}{\fontsize{9pt}{11pt}\selectfont}

% Amélioration des en-têtes et pieds de page
\usepackage{fancyhdr}
\fancypagestyle{thesis}{%
  \fancyhf{} % Clear all headers and footers
  \fancyhead[L]{\small\sffamily\leftmark}
  \fancyhead[R]{\small\sffamily\rightmark}
  \fancyfoot[C]{\sffamily\thepage}
  \renewcommand{\headrulewidth}{0.4pt}
  \renewcommand{\footrulewidth}{0pt}
}

% Style pour les tableaux amélioré
\usepackage{booktabs}
\usepackage{array}
\renewcommand{\arraystretch}{1.2} % Espacement vertical dans les tableaux

% Amélioration des titres de niveau inférieur
\titleformat{\paragraph}
  {\fontsize{11pt}{13pt}\selectfont\bfseries\color{blue!50!black}\sffamily}
  {\theparagraph}{1em}{}

\titleformat{\subparagraph}
  {\fontsize{10pt}{12pt}\selectfont\bfseries\color{blue!40!black}\sffamily}
  {\thesubparagraph}{1em}{}

% Style pour les environnements de code avec de meilleures fontes
\usepackage{listings}
\lstset{
    basicstyle=\footnotesize\ttfamily,
    commentstyle=\color{gray}\itshape,
    keywordstyle=\color{blue!80!black}\bfseries,
    stringstyle=\color{green!60!black},
    numberstyle=\tiny\color{gray},
    showstringspaces=false,
    breaklines=true,
    frame=single,
    framerule=0.5pt,
    rulecolor=\color{gray!50},
    backgroundcolor=\color{gray!5},
    tabsize=4,
    captionpos=b,
    xleftmargin=2em,
    xrightmargin=1em,
    framextopmargin=2pt,
    framexbottommargin=2pt
}
